\documentclass[a4paper,11pt]{article}

\usepackage[backend=biber,style=numeric-comp]{biblatex}
\addbibresource{ref.bib}

\usepackage[utf8]{inputenc}
\usepackage[T1]{fontenc}
\usepackage[french]{babel}

\usepackage[a4paper,
landscape,
top=2cm, bottom=2.5cm,
left=1.5cm, right=1.5cm]{geometry}

\usepackage{lmodern}
\usepackage{microtype}
\usepackage{booktabs}
\usepackage{longtable}
\usepackage{array}
\usepackage{xcolor}
\usepackage{colortbl}

\definecolor{headerblue}{HTML}{1F4E79}
\definecolor{rowgray}{HTML}{F2F2F2}

% \sp renommé en \espece pour éviter le conflit avec la commande LaTeX native
\newcommand{\espece}[1]{\textit{#1}}

\title{\textbf{Distribution verticale des larves et localisation des zones de ponte}\\[4pt]
	\large Synthèse bibliographique sur les 11 espèces cibles}
\author{}
\date{}

\begin{document}
	\maketitle
	\thispagestyle{empty}
	
	\setlength{\tabcolsep}{6pt}
	\renewcommand{\arraystretch}{1.5}
	
	\noindent\small
	\textbf{Abréviations :} EEC~= Eastern English Channel ;
	STST~= Selective Tidal Stream Transport ;
	DVM~= Diel Vertical Migration.
	
	\begin{longtable}{%
			>{\raggedright\arraybackslash}m{3.6cm}
			>{\raggedright\arraybackslash}m{5.2cm}
			>{\raggedright\arraybackslash}m{4.0cm}
			>{\raggedright\arraybackslash}m{8.0cm}
			>{\centering\arraybackslash}m{2.6cm}
		}
		
		\toprule
		\rowcolor{headerblue}
		\textcolor{white}{\textbf{Espèce}} &
		\textcolor{white}{\textbf{Profondeur des larves (m)}} &
		\textcolor{white}{\textbf{Profondeur de ponte (m)}} &
		\textcolor{white}{\textbf{Localisation des zones de ponte}} &
		\textcolor{white}{\textbf{Réf.}} \\
		\midrule
		\endfirsthead
		
		\toprule
		\rowcolor{headerblue}
		\textcolor{white}{\textbf{Espèce}} &
		\textcolor{white}{\textbf{Profondeur des larves (m)}} &
		\textcolor{white}{\textbf{Profondeur de ponte (m)}} &
		\textcolor{white}{\textbf{Localisation des zones de ponte}} &
		\textcolor{white}{\textbf{Réf.}} \\
		\midrule
		\endhead
		
		\midrule
		\endfoot
		
		\bottomrule
		\endlastfoot
		
		\rowcolor{white}
		\espece{Clupea harengus} &
		0--40 (jour) ; homogène ou $>$60 (nuit) &
		Benthique (non quantifiée) &
		Manche, sud Mer du Nord, estuaires &
		\textsuperscript{\cite{haslob_small_2009,hoffle_differences_2013,neven_winter_2024,ices_results_2007}} \\
		
		\rowcolor{rowgray}
		\espece{Gadus morhua} &
		0--40 (62\,\% en 0--20 de jour) &
		Épipelagique (non quantifiée) &
		Ensemble de la Mer du Nord &
		\textsuperscript{\cite{hoffle_differences_2013,ices_results_2007}} \\
		
		\rowcolor{white}
		\espece{Limanda limanda} &
		$>$40 (plus profond que les gadidés) &
		Pélagique (non quantifiée) &
		Sud-Est de la Mer du Nord &
		\textsuperscript{\cite{hoffle_differences_2013,ices_results_2007}} \\
		
		\rowcolor{rowgray}
		\espece{Melanogrammus aeglefinus} &
		0--40 (52\,\% en 0--20 de jour) &
		Pélagique (non quantifiée) &
		Nord Mer du Nord, côtes écossaises &
		\textsuperscript{\cite{hoffle_differences_2013,ices_results_2007}} \\
		
		\rowcolor{white}
		\espece{Merlangius merlangus} &
		0--40 (remontée 0--20 la nuit) &
		Pélagique (non quantifiée) &
		Large Mer du Nord, côtes anglaises &
		\textsuperscript{\cite{hoffle_differences_2013,ices_results_2007}} \\
		
		\rowcolor{rowgray}
		\espece{Microstomus kitt} &
		Pélagique (stades précoces) &
		16--92 (adultes capturés) &
		Manche Ouest, eaux écossaises &
		\textsuperscript{\cite{jennings_distribution_1993,ices_results_2007}} \\
		
		\rowcolor{white}
		\espece{Pleuronectes platessa} &
		Toute la colonne d'eau ; STST : fond au jusant, surface au flot &
		Épipelagique au large (non quantifiée) &
		Sud Mer du Nord, EEC, German Bight &
		\textsuperscript{\cite{de_graaf_numerical_2004,neven_winter_2024,ices_results_2007}} \\
		
		\rowcolor{rowgray}
		\espece{Pollachius virens} &
		$>$40 (jour) ; 0--40 (nuit) &
		Pélagique (non quantifiée) &
		Eaux profondes, nord Mer du Nord &
		\textsuperscript{\cite{hoffle_differences_2013,ices_results_2007}} \\
		
		\rowcolor{white}
		\espece{Scomber scombrus} &
		0--60 ($<$10\,mm) ; DVM marquée ($>$10\,mm) &
		0--30 (90\,\% des \oe{}ufs $<$13\,m) &
		Talus continental, Mer du Nord, Mer Celtique, Golfe de Gascogne &
		\textsuperscript{\cite{jansen_population_2013,jansen_temperature_2011,bartsch_individualbased_2004}} \\
		
		\rowcolor{rowgray}
		\espece{Sprattus sprattus} &
		Couches de surface &
		Pélagique (non quantifiée) &
		Sud Mer du Nord, German Bight &
		\textsuperscript{\cite{ices_results_2007}} \\
		
		\rowcolor{white}
		\espece{Trisopterus esmarkii} &
		$>$40 (jour) ; 75--100 (nuit) &
		Pélagique (non quantifiée) &
		Nord Mer du Nord, îles Shetland &
		\textsuperscript{\cite{hoffle_differences_2013,ices_results_2007}} \\
		
	\end{longtable}
	
	\vspace{0.8em}
	\noindent\small
	\printbibliography[title={Références}]
	
\end{document}